\section{Conclusion}


Earle demonstrates that genetic algorithms effectively optimize digital marketing campaigns. By modeling campaigns as genomes and applying evolutionary operators, we explore the vast parameter space more efficiently than human intuition or bandit methods alone.

Production deployments show 34\% conversion improvement and 2.8$\times$ faster convergence. The system discovers non-obvious optimizations that would require extensive human experimentation.

Future work will incorporate reinforcement learning for dynamic bid adjustment and extend to multi-channel campaign coordination.

\begin{thebibliography}{9}

\bibitem{holland1975}
J. H. Holland, \textit{Adaptation in Natural and Artificial Systems}, University of Michigan Press, 1975.

\bibitem{thompson1933}
W. R. Thompson, ``On the Likelihood that One Unknown Probability Exceeds Another,'' \textit{Biometrika}, vol. 25, pp. 285--294, 1933.

\bibitem{wordnet1995}
G. A. Miller, ``WordNet: A Lexical Database for English,'' \textit{Communications of the ACM}, vol. 38, no. 11, pp. 39--41, 1995.

\bibitem{xiao2014}
Y. Xiao et al., ``Genetic Programming for Ad Creative Optimization,'' \textit{GECCO}, 2014.

\bibitem{kumar2015}
R. Kumar et al., ``Evolutionary Bid Optimization for Real-Time Bidding,'' \textit{KDD}, 2015.

\bibitem{automl2019}
F. Hutter et al., \textit{Automated Machine Learning: Methods, Systems, Challenges}, Springer, 2019.

\end{thebibliography}

